\chapter{Plan de trabajo de la Etapa 1}
\label{ch:wbs_stage1}

\section{Alcance y criterio rector de la Etapa 1}

El plan de esta etapa se conserva como marco de actividades y entregables. Su avance científico actual se documenta en el \cref{ch:avances}: el banco dc01, B0--B2 y S4.1--S4.6d.2 desarrollan LT1 y el marco de evaluación de LT2, incluida la validación regional del predictor de retardo. Estos avances no implican que se hayan completado la instrumentación multimodal, el estudio con participantes o los hitos de vinculación comunitaria.

La Etapa 1 comprende los primeros diez meses del proyecto y corresponde a la fase de
\emph{modelado, diseño experimental y vinculación comunitaria}. Su propósito es
establecer las bases conceptuales, experimentales, metodológicas y éticas necesarias para
estudiar la recuperación de información fisiológica y de comportamiento a partir de
señales inalámbricas en escenarios reales. En consecuencia, la etapa articula seis tareas:
(i) simular escenarios Wi-Fi y LoRa; (ii) caracterizar las limitaciones físicas e
instrumentales; (iii) definir el conjunto de datos multimodal; (iv) formalizar el protocolo
de adquisición; (v) establecer métodos de referencia con simulación, datos abiertos y
pruebas preliminares; y (vi) seleccionar y caracterizar el sitio piloto mediante actividades
de vinculación.

Antes del despliegue comunitario se realizará una campaña experimental controlada en
laboratorio. Esta campaña constituye el vínculo empírico entre la simulación y el piloto:
permite verificar la instrumentación, la sincronización, la repetibilidad, la observabilidad
y los procedimientos de calidad. No sustituye la adquisición formal prevista para la
Etapa 2 ni pretende anticipar sus resultados; reduce la incertidumbre técnica antes de
trasladar el protocolo a una Colmena u otro espacio comunitario de Guadalajara.

El éxito de la Etapa 1 no se juzgará por el volumen de datos ni por el número de
participantes, sino por la reducción verificable de la incertidumbre metodológica. Al
finalizar el mes 10 deberán estar disponibles un modelo físico reproducible, una plataforma
de adquisición estable, un estudio preliminar de laboratorio documentado, un protocolo
validado y un sitio comunitario caracterizado para la campaña formal de la Etapa 2. Esta
secuencia mantiene la continuidad con la metodología general: primero se determina qué
debería ser observable, después qué puede medirse y, por último, bajo qué condiciones es
justificado avanzar al entorno comunitario.

\section{Organización de las líneas de trabajo}

La Etapa 1 se organiza en seis líneas de trabajo (LT). La organización conserva los
objetivos E1-1--E1-6 del proyecto y los agrupa de acuerdo con sus dependencias científicas
y operativas. Las líneas no son compartimentos independientes: el modelado fundamenta
el protocolo; el protocolo determina la instrumentación; la instrumentación hace posible
el estudio preliminar; y la evidencia reunida condiciona la transición a la Etapa 2.

\begin{table}[H]
\centering
\caption{Líneas de trabajo de la Etapa 1 y correspondencia con los objetivos específicos.}
\label{tab:wbs_stage1_summary}
\small
\begin{tabularx}{\textwidth}{P{1.15cm}P{3.25cm}P{1.9cm}L}
\toprule
\textbf{Línea} & \textbf{Denominación} & \textbf{Objetivos} & \textbf{Resultado esperado}\\
\midrule
LT1 & Modelado físico y gemelo inalámbrico inicial
& E1-1, E1-2
& Entorno reproducible de simulación Wi-Fi/LoRa, comparación con mediciones reales y métricas de propagación, fase y observabilidad.\\

LT2 & Marco de datos, protocolo y métodos de referencia
& E1-3, E1-4, E1-5
& Especificación multimodal, protocolo v1.0, esquema de metadatos, control de calidad y métodos reproducibles de referencia.\\

LT3 & Plataforma experimental multimodal de laboratorio
& E1-2, E1-4
& Banco experimental integrado y caracterizado: Wi-Fi, LoRa, referencias mmWave/visión/fisiología y sincronización.\\

LT4 & Estudio experimental preliminar y controlado
& E1-2, E1-5
& Conjunto preliminar de laboratorio y validación de presencia, movimiento, micromovimiento y respiración en condiciones controladas.\\

LT5 & Vinculación comunitaria, Colmenas y taller de IoT
& E1-6
& Sitio piloto caracterizado, caso de uso comunitario seleccionado y taller práctico ejecutado en Guadalajara.\\

LT6 & Integración, evaluación y transición a la Etapa 2
& E1-3--E1-6
& Informe técnico de Etapa 1, paquete de transición, matriz de riesgos y protocolo para despliegue formal.\\
\bottomrule
\end{tabularx}
\end{table}

\subsection{LT1 --- Modelado físico y gemelo inalámbrico inicial}
\label{sec:wp1_stage1}

\textbf{Periodo:} M1--M6; cierre inicial S4.7a en M1--M2 y ampliación geométrica en M4--M6.

\textbf{Avance documentado:} S1 y B0--B2 cuantifican la discrepancia estática; S4.1--S4.6d.2 analizan fase, estructura residual y transferencia del retardo al canal. S4.6 caracteriza la fase común, recupera el tiempo de adquisición y evalúa continuidad local. El trabajo siguiente S4.7a estudiará fidelidad relativa; S4.7b abordará incertidumbre por trayectoria. Esto aporta evidencia a P1.2 y P1.3. P1.1 y la revisión R1 requieren vincular los registros reproducibles indicados en la \cref{sec:trazabilidad_avances}.

\textbf{Pregunta rectora:} ¿qué información sobre entorno, movimiento y micromovimiento
puede preservar el canal Wi-Fi/LoRa bajo geometrías y configuraciones de radio
controladas?

\paragraph{Actividades.}
\begin{enumerate}[label=\textbf{LT1.\arabic*:},leftmargin=3.8em]
    \item Construcción y validación de casos de referencia: espacio libre, reflexión,
    multitrayectoria y perturbaciones geométricas controladas.
    \item Reproducción del banco de pruebas DICHASUS dc01 y generación del conjunto
    $(\vect p_i,H_i^{\rm real},\mathcal P_i^{\rm RT})$.
    \item Evaluación del trazado de rayos nominal y calibración progresiva de materiales,
    dispersión y errores locales de fase.
    \item Implementación del modelo de observación Wi-Fi OFDM y del modelo inicial
    LoRa CSS/IQ.
    \item Definición de métricas de fidelidad de campo y sensibilidad diferencial,
    incluyendo NMSE complejo, error de fase, PDP, dispersión cuadrática media del retardo y
    $\partial H/\partial q$.
    \item Construcción de mapas preliminares de cobertura y observabilidad en
    escenarios de referencia en interiores y exteriores.
\end{enumerate}

\paragraph{Productos y criterio de aceptación.}
\begin{itemize}
    \item \textbf{P1.1} Repositorio reproducible del modelo directo.
    \item \textbf{P1.2} Comparación DICHASUS entre mediciones y simulación.
    \item \textbf{P1.3} Documento técnico de limitaciones físicas e instrumentales.
    \item \textbf{Revisión R1:} las simulaciones son reproducibles y la discrepancia
    entre simulación y medición se encuentra cuantificada mediante métricas definidas de antemano.
\end{itemize}

\subsection{LT2 --- Marco de datos, protocolo y métodos de referencia}
\label{sec:wp2_stage1}

\textbf{Periodo:} M1--M5.

\textbf{Pregunta rectora:} ¿qué información debe conservarse y bajo qué reglas debe
adquirirse para que las mediciones sean comparables, trazables y reutilizables?

\paragraph{Actividades.}
\begin{enumerate}[label=\textbf{LT2.\arabic*:},leftmargin=3.8em]
    \item Revisión y reproducción de flujos de procesamiento sobre conjuntos de datos públicos de canal,
    actividad y fisiología.
    \item Definición del modelo de datos: señal de radiofrecuencia, valores de referencia, estado del participante,
    contexto, geometría, configuración de radio y metadatos.
    \item Especificación de nomenclatura, control de versiones, marcas temporales, estructura de
    archivos, trazabilidad y criterios de inclusión/exclusión.
    \item Diseño del protocolo multimodal para configuraciones sin participantes y
    con uno o varios participantes.
    \item Definición de procedimientos de sincronización y control de calidad.
    \item Formalización de consentimiento, anonimización, gestión de video,
    acceso y gobernanza de datos.
    \item Implementación de métodos de referencia sin selección oráculo de canales,
    subportadoras o componentes.
\end{enumerate}

\paragraph{Productos y criterio de aceptación.}
\begin{itemize}
    \item \textbf{P2.1} Diccionario y esquema del conjunto de datos multimodal.
    \item \textbf{P2.2} Protocolo de adquisición v1.0.
    \item \textbf{P2.3} Lineamientos de ética, privacidad y gobernanza.
    \item \textbf{P2.4} Métodos reproducibles de referencia sobre datos públicos.
    \item \textbf{Revisión R2:} no se iniciará una campaña formal con participantes sin
    protocolo, trazabilidad, control de calidad y autorización ética aplicable.
\end{itemize}

\subsection{LT3 --- Plataforma experimental multimodal de laboratorio}
\label{sec:wp3_stage1}

\textbf{Periodo:} M1--M6; diseño y adquisición M1--M3, caracterización E0 en M2--M4.

\textbf{Objetivo:} implementar una plataforma común que permita observar el mismo
fenómeno físico mediante radios de interés y modalidades de referencia.

La arquitectura experimental se define como

\begin{equation}
\begin{aligned}
\mathcal A_{\rm lab}=\{&Y_{\rm WiFi},Y_{\rm LoRa},Y_{\rm mmWave},Y_{\rm ref},\\
&V_{\rm RGB(D)},x_{\rm stage},x_{\rm resp},x_{\rm ECG/RR},x_{\rm IMU}\}.
\end{aligned}
\label{eq:lab_architecture}
\end{equation}

Esta arquitectura es un catálogo de modalidades posibles; cada captura declarará el
subconjunto realmente disponible. $Y_{\rm ref}$ y $x_{\rm stage}$ corresponden a la
referencia RF y al desplazamiento mecánico medido. ECG e inerciales son auxiliares
según tarea y no requisitos para cerrar el banco mecánico.

\paragraph{Bloques instrumentales.}
\begin{description}[style=nextline]
    \item[\textbf{Referencia RF y mecánica.}]
    Adquisición coherente caracterizada, etapa con medición independiente y soportes
    registrados. La incertidumbre del banco determina los pasos utilizables en E1.
    \item[\textbf{Wi-Fi.}]
    Enlaces configurables para CSI, con control de canal, ancho de banda,
    geometría Tx--Rx, orientación y tasa de paquetes.

    \item[\textbf{LoRa/LoRaWAN.}]
    Nodos y puerta de enlace de laboratorio; cuando sea posible, observación I/Q mediante
    SDR para preservar información anterior a la reducción a RSSI/SNR.

    \item[\textbf{Radar de ondas milimétricas.}]
    Modalidad de radiofrecuencia de referencia para movimiento y micromovimiento.

    \item[\textbf{RGB/RGB-D.}]
    Referencia geométrica para posición, orientación, trayectoria y actividad.

    \item[\textbf{Fisiología e inerciales.}]
    Banda de expansión respiratoria, ECG o banda torácica con intervalos RR,
    y una o más IMU para movimiento corporal.
\end{description}

\paragraph{Actividades de caracterización.}
Se medirán estabilidad temporal, piso de ruido, pérdida de paquetes, CFO/SFO/CPO,
latencia entre modalidades y deriva temporal. La sincronización se expresará como

\begin{equation}
a_m t_i^{(m)}+b_m=t_i^{\rm global}+\epsilon_i^{(m)},
\end{equation}

donde $a_m,b_m$ corrigen escala y origen del reloj. El error residual
$\epsilon_i^{(m)}$ deberá estimarse y reportarse, no asumirse nulo; se segmentarán
las sesiones si los reinicios o la deriva invalidan un único mapa afín.

\paragraph{Productos y criterio de aceptación.}
\begin{itemize}
    \item \textbf{P3.1} Banco experimental inventariado y documentado.
    \item \textbf{P3.2} Procedimiento de calibración y sincronización.
    \item \textbf{P3.3} Informe de estabilidad y repetibilidad.
    \item \textbf{Revisión R3:} las modalidades del banco mínimo E0--E1 producen
    registros trazables, con referencia y sincronización caracterizadas. La integración
    adicional continuará según el cronograma y se validará antes de utilizarla.
\end{itemize}

\subsection{LT4 --- Estudio experimental preliminar y controlado}
\label{sec:wp4_stage1}

\textbf{Periodo:} M3--M9; reflector M3--M5, fantoma M5--M7 y una persona M7--M9.

El estudio preliminar se ejecutará de forma acumulativa. Cada nivel habilita el siguiente; no se
introducirán participantes antes de validar primero el comportamiento instrumental y
geométrico.

\begin{table}[H]
\centering
\caption{Matriz experimental mínima del estudio preliminar de laboratorio.}
\label{tab:lab_prepilot}
\small
\begin{tabularx}{\textwidth}{P{1.1cm}P{3.0cm}L L}
\toprule
\textbf{Nivel} & \textbf{Condición} & \textbf{Variable controlada} & \textbf{Resultado buscado}\\
\midrule
N0 & Habitación vacía & Tiempo, canal, dispositivo & Piso de ruido, estabilidad y perturbaciones instrumentales.\\
N1 & Reflector estático & Posición y material & Sensibilidad espacial y repetibilidad.\\
N2 & Reflector móvil & Desplazamiento conocido & Fidelidad diferencial y $\partial H/\partial x$.\\
N3 & Fantoma respiratorio & Amplitud, frecuencia y movimiento lento & Forma de onda mecánica y sensibilidad reproducibles.\\
N4 & Una persona & Postura, actividad y respiración & Movimiento y forma de onda frente a referencias sincronizadas.\\
N5 & Wi-Fi y LoRa pareados & Misma perturbación física & Comparación de sensibilidad y observabilidad entre tecnologías.\\
N6 & Dos personas, exploratorio & Separación y actividad & Viabilidad y condicionamiento antes de la Etapa 2.\\
\bottomrule
\end{tabularx}
\end{table}

El conjunto de salida se define como

\begin{equation}
\begin{aligned}
\mathcal D_{\rm lab}^{(0)}=\{&Y_{\rm WiFi},Y_{\rm LoRa},Y_{\rm mmWave},Y_{\rm ref},\\
&V_{\rm RGB(D)},x_{\rm stage},x_{\rm resp},x_{\rm ECG/RR},x_{\rm IMU},M\},
\end{aligned}
\end{equation}

donde $M$ contiene metadatos y estado de control de calidad.

\paragraph{Revisión R4.}
En M7 se revisará la evidencia de E0--E2: estabilidad instrumental, respuesta a
movimiento medido y recuperación del fantoma con incertidumbre. La adquisición de
una persona se prevé en M7--M9 y requiere además aprobación ética. El cumplimiento
del hito mecánico no equivale a validación fisiológica ni habilita automáticamente
un despliegue humano comunitario.

\subsection{LT5 --- Vinculación comunitaria, Colmenas Guadalajara y taller de IoT}
\label{sec:wp5_stage1}

\textbf{Periodo:} M6--M10; preparación institucional desde M6 y taller previsto en M9--M10.

\textbf{Objetivo:} seleccionar y caracterizar un sitio comunitario en Guadalajara y
establecer una relación de trabajo que combine preparación del piloto con apropiación
tecnológica.

La vinculación seguirá cuatro fases:
\begin{enumerate}[label=\textbf{V\arabic*:},leftmargin=3em]
    \item \textbf{Vinculación institucional:} identificación de contraparte local,
    condiciones de acceso, responsables y logística.
    \item \textbf{Caracterización técnica:} geometría, materiales, energía,
    infraestructura Wi-Fi, cobertura LoRa, ruido RF, zonas de actividad y restricciones
    de instalación.
    \item \textbf{Codiseño comunitario:} selección de un único caso de uso ambiental
    para el nodo IoT de divulgación.
    \item \textbf{Taller y preparación del despliegue:} taller práctico, demostración de sensores,
    conectividad y energía autónoma, y pruebas técnicas sin convertir el taller en una
    campaña experimental de sujetos humanos.
\end{enumerate}

El resultado técnico de la visita será una ficha de sitio y un primer modelo
$G_E^{\rm Colmena}$ utilizable posteriormente para simulación y planificación de la
Etapa 2.

\paragraph{Revisión R5.}
El sitio deberá contar con contraparte identificada, caracterización técnica, restricciones
operativas documentadas y un caso de uso comunitario definido.

\subsection{LT6 --- Integración y transición a la Etapa 2}
\label{sec:wp6_stage1}

\textbf{Periodo:} M8--M10.

La línea final consolida los resultados de las cinco líneas anteriores y establece la primera
versión operativa del sistema.

\paragraph{Productos de cierre.}
\begin{enumerate}[label=\textbf{P6.\arabic*:},leftmargin=3em]
    \item Entorno de simulación reproducible Wi-Fi/LoRa.
    \item Plataforma multimodal de laboratorio caracterizada.
    \item Protocolo de adquisición, metadatos y control de calidad v1.0.
    \item Conjunto preliminar $\mathcal D_{\rm lab}^{(0)}$.
    \item Métodos reproducibles de referencia.
    \item Sitio comunitario caracterizado.
    \item Taller y materiales de vinculación ejecutados.
    \item Matriz de riesgos y pendientes.
    \item Plan cerrado de despliegue y adquisición para la Etapa 2.
\end{enumerate}

La transición se formaliza mediante

\begin{equation}
R_{\rm E2}=
R_{\rm RF}\land
R_{\rm sinc}\land
R_{\rm protocolo}\land
R_{\rm etica}\land
R_{\rm sitio}\land
R_{\rm datos}.
\label{eq:stage2_gate}
\end{equation}

La Etapa 2 deberá iniciar su campaña formal cuando estas condiciones hayan sido
verificadas, y no únicamente por cumplimiento cronológico del mes 10.

\section{Secuencia y dependencias}

Las dependencias principales siguen una secuencia de evidencia. LT1 establece las
predicciones físicas que deben contrastarse; LT2 determina cómo adquirir y conservar las
mediciones; LT3 materializa esas decisiones en una plataforma común; y LT4 contrasta el
modelo con perturbaciones controladas. LT5 se inicia cuando existe suficiente madurez
técnica para caracterizar el sitio comunitario, mientras que LT6 integra la evidencia y
documenta los supuestos pendientes. El solapamiento temporal mostrado en la
\cref{tab:cronograma_etapa1} expresa trabajo concurrente, no independencia entre líneas.

\section{Cronograma de la Etapa 1}

\subsection{Paquete técnico de las primeras ocho semanas}
\label{sec:eight_week_package}

Dentro de M1--M2 se propone un paquete corto que convierte la revisión de literatura
en artefactos verificables antes de cerrar especificaciones de compra:
\begin{table}[htbp]
\centering\small
\caption{Secuencia técnica inmediata; las semanas son relativas al inicio del paquete.}
\label{tab:eight_week_package}
\begin{tabularx}{\textwidth}{P{1.5cm}L L}
\toprule
Semana & Actividad & Evidencia de cierre\\
\midrule
1 & Congelar versiones, manifiestos, particiones y artefactos S1--S4.6. &
Correspondencia reproducible y sumas de verificación.\\
2 & Implementar $\mathsf V_0$--$\mathsf V_7$ y $\mathsf T_1$--$\mathsf T_7$. &
Pruebas algebraicas y matriz de respuesta a perturbaciones.\\
3 & Ejecutar fidelidad estática de dc01. & Resultados densos y por bloques con
normalización ajustada sólo en entrenamiento.\\
4 & Auditar d010 y realizar pruebas de sensibilidad si sus metadatos lo permiten. &
Pareto entre rechazo de perturbaciones y cambio geométrico del transmisor.\\
5 & Implementar S4.7b. & Baselines circulares, NLL, calibración y abstención.\\
6 & Integrar perturbaciones seleccionadas de RadioRange. & Vulnerabilidades por
observable bajo rangos trazables.\\
7 & Ejecutar un caso mínimo de Indoor Radio Mapping y WiSegRT. & Flujo de nube/malla
a RT; RSSI y canal simulado con límites declarados.\\
8 & Auditar subconjuntos CAEZ y OpenLoRa. & Adaptadores iniciales y requisitos
medidos para el banco propio.\\
\bottomrule
\end{tabularx}
\end{table}

El paquete no condiciona toda adquisición a obtener un método ganador. Debe responder
qué fase necesita referencia, qué ancho de banda y coherencia preservan sensibilidad,
qué precisión geométrica resulta relevante y qué salidas proporciona cada radio.
Las compras de baja ambigüedad necesarias para mecánica, montaje y trazabilidad pueden
prepararse en paralelo; las especificaciones de recepción coherente se cerrarán con
la evidencia disponible y las cotizaciones compatibles.

\begin{table}[htbp]
\centering\small
\caption{Cronograma revisado: meses relativos al inicio operativo de la etapa; ventanas propuestas sujetas a evidencia.}
\label{tab:cronograma_etapa1}
\begin{tabularx}{\textwidth}{P{1.8cm}P{2.2cm}L L}
\toprule
Meses & Líneas & Actividad & Producto revisable\\
\midrule
M1--M2 & LT1/LT2 & S4.7a en dc01. & Comparación $\mathsf V_i$ y controles espaciales.\\
M1--M3 & LT2/LT3 & Diseño del banco y compras prioritarias. & Referencia RF, metrología y captura especificadas.\\
M2--M4 & LT3 & Esc. E0. & Deriva y repetibilidad documentadas.\\
M3--M5 & LT3/LT4 & Esc. E1, reflector. & $\Delta H$ y Jacobiano con desplazamiento medido.\\
M4--M6 & LT1/LT3 & Reconstrucción RGB-D/LiDAR. & G0--G4, registro y controles independientes.\\
M5--M7 & LT4 & Esc. E2, fantoma. & Fidelidad diferencial y forma de onda.\\
M6--M8 & LT1/LT4 & Wi-Fi, LoRa y radar comparativos. & Operadores y referencias bajo perturbación común.\\
M7--M9 & LT4 & Esc. E3 con aprobación ética. & Datos humanos preliminares y límites.\\
M8--M10 & LT6 & Integración y transferencia inicial E5. & Protocolo y conjunto v1; pendientes de Etapa 2.\\
M6--M10 & LT5 & Vinculación; taller en M9--M10. & Sitio caracterizado y demostración de divulgación.\\
\bottomrule
\end{tabularx}
\end{table}

El calendario expresa dependencias técnicas y no acredita ejecución por transcurso
del tiempo. Esc. E4 y transferencia externa amplia se incorporarán según madurez y
recursos; no se prometen como resultados humanos completos dentro de esta ventana.
Mientras se integra el equipo se puede avanzar con CAEZ, el flujo geometría--RSSI,
OpenLoRa, reconstrucción RGB-D, marcas temporales y diseño de referencia. Cada conjunto
adicional se someterá primero a una auditoría de observables y procedencia.

\section{Movilidad científica}

El presupuesto aprobado considera una bolsa de \$68,000 MXN para viáticos y
\$30,000 MXN para pasajes nacionales, originalmente vinculados a actividades de
divulgación y vinculación. La planeación operativa propone concentrar estas movilidades
en una ventana M9--M10, de modo que una misma visita contribuya a: (i) revisión
experimental; (ii) validación Wi-Fi/LoRa; (iii) caracterización de la Colmena;
(iv) impartición del taller; y (v) preparación de la Etapa 2.

La solicitud aprobada refiere hospedaje y alimentación para dos personas por 12 noches y
tres vuelos nacionales redondos. Por tanto, una estancia presencial completa de un mes
para el especialista deberá considerarse una extensión deseable, pero su ejecución requerirá
validar previamente con la administración del proyecto la posibilidad de redistribución del
rubro o complementarla con recursos institucionales. El documento operativo no debe
suponer dicha ampliación como autorizada.

La movilidad prospectiva a UPNA constituye una ampliación científica distinta de los
pasajes nacionales presupuestados. Requerirá acuerdo con el grupo anfitrión, disponibilidad
de instrumentación y financiación específica; su objetivo será repetir E1/E2 con el
protocolo ya validado localmente. Una propuesta de colaboración no se registra como
estancia concedida ni como validación realizada.

\section{Plan de adquisiciones}

El presupuesto aprobado contiene tres bolsas materiales diferenciadas:
\$40,000 MXN para dispositivos electrónicos de radio configurables y accesorios;
\$40,000 MXN para sensores y suministros de laboratorio de referencia; y
\$70,000 MXN para 25 kits electrónicos de divulgación.


\subsection{Inventario reportado y verificación funcional}

La planeación considera como inventario reportado plataformas Wi-Fi/ESP32, nodos
LoRa STM32WL, equipos Notecard y RAK2287, un ADALM-Pluto, tres IWRL6432BOOST y
cámaras RGB, RGB-D o estéreo. Antes de comprar se registrarán modelos exactos,
interfaces y estado funcional. La denominación Notecard no basta para identificar
tecnología o acceso a señal; se comprobará el SKU. Tampoco se asumirá que un
concentrador LoRaWAN entrega I/Q coherente.

Los IWRL6432BOOST se basan en la familia radar de 57--64 GHz
\citep{TIIWRL6432BOOST}. Su disponibilidad reportada permite plantear vistas complementarias;
se verificarán firmware, acceso a datos, calibración y sincronización antes de fijar
el protocolo. Su precisión de desplazamiento deberá medirse y no se adoptará como
referencia submilimétrica independiente. Se reutilizarán estas unidades antes de
priorizar otro radar equivalente.

\subsection{Plataforma experimental de radiofrecuencia --- bolsa autorizada: \$40,000 MXN}

\begin{table}[htbp]
\centering\small
\caption{Prioridades RF revisadas a partir del inventario reportado.}
\label{tab:rf_procurement}
\begin{tabularx}{\textwidth}{P{3.4cm}L P{2.9cm}}
\toprule
Bloque & Función y condición & Prioridad\\
\midrule
Referencia coherente & Resolver recepción simultánea o referencia verificable; comprobar
compatibilidad con SDR y reloj antes de seleccionar equipo. & Muy alta\\
Coaxiales, divisores, atenuadores y adaptadores & Caracterizar pérdidas/fase y estabilizar
cadenas; cantidades según arquitectura, sin asumir calibración de fábrica suficiente. & Muy alta\\
Antenas y soportes rígidos & Registro geométrico, polarización y repetibilidad del centro
de montaje; cables inmovilizados. & Alta\\
Sincronización y automatización & Marcas temporales, disparos disponibles y control
mecánico; medición de deriva entre modalidades. & Alta\\
Ampliación Wi-Fi/LoRa/SDR & Cubrir una carencia demostrada de CSI, I/Q o enlaces;
reutilizar primero plataformas existentes. & Condicionada\\
\bottomrule
\end{tabularx}
\end{table}

La referencia es un requisito funcional, no una compra cerrada: el Pluto documentado
no ofrece por sí solo dos receptores simultáneos. Se compararán soluciones compatibles
y su coste total dentro del rubro. Compartir frecuencia o PPS no exime de verificar
fase relativa tras reinicios.

\subsection{Metrología y referencias --- bolsa autorizada: \$40,000 MXN}

\begin{table}[htbp]
\centering\small
\caption{Prioridades para metrología y referencias del banco actualizado.}
\label{tab:reference_procurement}
\begin{tabularx}{\textwidth}{P{3.4cm}L P{2.9cm}}
\toprule
Instrumento & Uso científico & Prioridad\\
\midrule
Etapa lineal y referencia micrométrica & Desplazamiento medido, control de holgura,
reflector y fantoma; incertidumbre compatible con E1. & Muy alta\\
Fijaciones y marcadores & Registro reproducible de escena, sensores y antenas;
AprilTags u otros fiduciales con dimensiones verificadas. & Muy alta\\
Referencia respiratoria & Forma de onda sincronizada para E3; no sustituible por una
etiqueta de BPM o ECG aislado. & Alta\\
LiDAR 3D de gama media & Geometría global y repetibilidad multiescena; comparar G1--G3
frente a RGB-D existente. & Alta; cotizar\\
Distanciómetro o control métrico & Comprobación independiente de dimensiones; Bosch GLM
o equivalente es opcional, no requisito de marca. & Complementaria\\
Radar, RGB-D, IMU y ECG & Reutilizar disponibilidad; ampliar sólo ante una necesidad
medida. Cardiología queda fuera del núcleo E0--E2. & Condicionada\\
\bottomrule
\end{tabularx}
\end{table}

Para LiDAR se propone evaluar una clase Livox Mid-360 o equivalente, sin fijar compra,
precio ni exactitud no comprobados. La especificación exigirá exportación de nube de
puntos, marcas temporales accesibles, integración mediante SDK/ROS2 o interfaz compatible,
registro métrico y documentación de incertidumbre. No se requiere un equipo topográfico
ni se presupone que un escáner de mayor coste resolverá la fase del canal.

La selección monetaria se realizará contra cotizaciones y compatibilidad real. La
combinación completa de instrumentos no se da por financiada dentro de \$40,000 MXN.
Si excede el rubro, se priorizarán mecánica, referencia y sincronización, usando RGB-D
para iniciar el registro y gestionando préstamo o recursos complementarios para LiDAR.
No se trasladarán recursos entre bolsas ni se imputará movilidad internacional a pasajes
nacionales mediante esta actualización metodológica.

\subsection{Kits de divulgación --- bolsa autorizada: \$70,000 MXN}

El presupuesto aprobado corresponde a 25 kits, equivalente a un promedio de

\begin{equation}
\bar C_{\rm kit}=
\frac{70,000}{25}=2,800\ {\rm MXN/kit}.
\end{equation}

Cada kit se diseñará como plataforma modular con microcontrolador, conectividad,
sensores, alimentación y material de prototipado. El sensor específico del caso de uso
ambiental deberá seleccionarse después de la caracterización de la Colmena, evitando
adquirir anticipadamente 25 sensores especializados para un problema todavía no
seleccionado.

\section{Presupuesto consolidado de la Etapa 1}

\begin{table}[H]
\centering
\caption{Presupuesto aprobado y asignación por líneas de trabajo durante la Etapa 1.}
\label{tab:stage1_budget_formal}
\small
\begin{tabularx}{\textwidth}{L r P{2.1cm}}
\toprule
\textbf{Rubro} & \textbf{Monto MXN} & \textbf{Líneas principales}\\
\midrule
Apoyo a la investigación científica
& \$240,000 & LT1--LT6\\
Honorarios por servicios profesionales
& \$240,000 & LT1--LT6\\
Material eléctrico/electrónico para kits del taller
& \$70,000 & LT5\\
Viáticos
& \$68,000 & LT4--LT5\\
Pasajes nacionales
& \$30,000 & LT4--LT5\\
Material eléctrico/electrónico para pilotos
& \$40,000 & LT3--LT4\\
Sensores y suministros de laboratorio
& \$40,000 & LT3--LT4\\
\midrule
\textbf{Total Etapa 1}
& \textbf{\$728,000}
& \\
\bottomrule
\end{tabularx}
\end{table}

El cuadro anterior conserva íntegramente los rubros y montos aprobados. Las
subasignaciones internas de equipo se consideran planeación técnica y deberán cerrarse
contra cotizaciones y reglas administrativas antes de comprometer el gasto.

\section{Matriz de hitos y decisiones}

\begin{table}[H]
\centering
\caption{Hitos de decisión de la Etapa 1.}
\label{tab:stage1_milestones}
\small
\begin{tabularx}{\textwidth}{P{1.7cm}P{1.4cm}L L}
\toprule
\textbf{Revisión} & \textbf{Mes} & \textbf{Evidencia requerida} & \textbf{Decisión resultante}\\
\midrule
R1 & M4
& Gemelo estático y S4.7a reproducibles; discrepancia y acceso a fase declarados.
& Establecer el modelo físico de referencia.\\
R2 & M5
& Protocolo, metadatos, control de calidad y ética formalizados.
& Autorizar el estudio preliminar con participantes cuando corresponda.\\
R3 & M5
& Banco mínimo E0--E1 estable, referencia y sincronización caracterizadas.
& Continuar fantoma e integración multimodal, con control de deriva.\\
R4 & M7
& E0--E2: repetibilidad, sensibilidad mecánica y fantoma con incertidumbre.
& Iniciar E3 sólo con ética y referencias verificadas; preparar transferencia.\\
R5 & M9
& Colmena caracterizada, contraparte y caso de uso definidos.
& Preparar la instalación de la Etapa 2.\\
$R_{\rm E2}$ & M10
& Cumplimiento conjunto de radiofrecuencia, sincronización, protocolo, ética, sitio y datos.
& Iniciar campaña formal de Etapa 2.\\
\bottomrule
\end{tabularx}
\end{table}

\section{Síntesis}

La Etapa 1 constituye una fase de reducción de incertidumbre científica, no sólo de
cumplimiento administrativo. El modelado establece qué debería observarse; la
instrumentación verifica qué puede medirse; el estudio preliminar determina qué
información se conserva en condiciones reales; y la vinculación comunitaria procura que
el despliegue posterior ocurra en un sitio técnicamente caracterizado y con una contraparte
local. En conjunto, estas actividades convierten los fundamentos desarrollados en los
capítulos anteriores en una secuencia experimental auditable y preparan, sin anticiparla,
la validación de la Etapa 2.
